\section{Square products with short selections}
\label{p13:short:section}

The conditional count in the next section must handle a relation
even when none of its selected integers has a large prime kernel.
For selections of at most a small multiple of \(\log\log N\),
we instead count square products through moments. The constants
in this argument are explicit in the selection size; a
fixed-degree result alone would not give the required range.

The quadratic kernel argument is related to the decomposition
in de la Bret\`eche, Wang, and Xu
\cite[Lemmas~6.2--6.3]{bretwngxu2026}. We give the estimates
needed here, including their dependence on the number and
location of the roots.


Write
\[
 P_A(X)=\prod_{a\in A}(X+a),\qquad
 A\subset[-DB,DB]\cap\Z,
 \qquad 2\le d=|A|\le c\log\log N,
\]
where $d$ is even and $D,c>0$ are fixed.  The interval of summation
has height and length comparable to $N$.

\begin{proposition}[Uniform moments for short split products]
\label{p13:short:moments}
For every fixed even integer \(v\ge2\), with the preceding notation,
\begin{equation}
 \E\left|\sum_{n\asymp N}f(P_A(n))\right|^v
 \le N^{v/2+c[v(v-1)+2]\log2+o(1)}.
 \label{p13:short:moment-bound}
\end{equation}
In particular, for each fixed \(v\) and \(\varepsilon>0\),
one may choose \(c>0\) so that the \(L^v\) norm is
\(O(N^{1/2+\varepsilon})\), uniformly in all the indicated
motifs. The summation interval may have any fixed comparable
length and positive height.
\end{proposition}

The proof, in Section~\ref{p14:comp:short:moments}, pairs the linear
factors into quadratics and controls the common kernel by their
resultant.  In the moment expansion, primes are assigned to the even
subsets of indices on which they occur.  Averaging the exposure order
then gives the exponent $v/2$.  The explicit term
$c[v(v-1)+2]\log2$ is what permits a growing selection size.

\begin{corollary}[Short projections in an old space]
\label{p13:short:old-space}
Fix \(r,s,D\) and \(\varepsilon>0\). Let \(V_{\rm old}\)
be generated by the square classes of all values in \(s\)
fixed blocks of diameter and cardinality \(O_D(B)\).
For a sufficiently small fixed \(c>0\), the number of
lists of \(r\) new centres admitting nonempty even
selections \(A_i\), \(|A_i|\le c\log\log N\), whose
square classes sum to an element of \(V_{\rm old}\), is
\(O_{r,s,D,C,\varepsilon}(N^{r/2+\varepsilon})\),
uniformly in the old blocks.
\end{corollary}

\begin{proof}
For a square class \(u\), let \(\chi_u(f)\) denote its
character, and put \(\Phi_V=\sum_{u\in V}\chi_u\), summing
over distinct classes. Character orthogonality gives
\[
 \Phi_V\ge0,\qquad \E\Phi_V=1,\qquad
 \|\Phi_V\|_a=2^{\dim V(1-1/a)}\quad(a\ge1).
\]
For fixed motifs, with \(X_i=\sum_z f(P_{A_i}(z))\),
the expectation \(\E(\Phi_{V_{\rm old}}\prod_iX_i)\)
counts exactly the required lists. For each list, a unique
old character cancels the new sum if it belongs to the old
space. All terms after expectation are zero or one.

Since \(\dim V_{\rm old}=O_{s,D}(B)\), choose a fixed
\(a>1\) sufficiently close to one that
\(\|\Phi_{V_{\rm old}}\|_a\le N^{\varepsilon'}\).
Next choose a fixed even \(v\ge ra/(a-1)\), and then choose
\(c\) by Proposition~\ref{p13:short:moments}.
H\"older gives
\[
 \left|\E\left(\Phi_{V_{\rm old}}\prod_iX_i\right)\right|
 \le N^{\varepsilon'}\prod_i
       \|X_i\|_{ra/(a-1)}
 \ll N^{r/2+(r+1)\varepsilon'}.
\]
The motifs have
\(\exp(O_{r,D}(c(\log\log N)^2))=N^{o(1)}\) choices.
Choose the internal margins in the stated \(\varepsilon\).
There is no enumeration of the \(2^{\dim V_{\rm old}}\)
old classes.
\end{proof}
